\documentclass[12pt]{article}

\usepackage{sbc-template}
\usepackage{graphicx,url}
\usepackage[utf8]{inputenc}
\usepackage[brazil]{babel}
\usepackage{booktabs}
\usepackage{amsmath}

\sloppy

\title{Sistema de Recuperação de Informação Híbrido para\\Artigos sobre IA Aplicada a Compras Públicas}

\author{Luiz Fernando Postingel Quirino\inst{1}}

\address{Faculdade de Computação -- Universidade Federal de Mato Grosso do Sul (UFMS)\\
  Campo Grande -- MS -- Brasil
  \email{luiz.quirino@ufms.br}
}

\begin{document}

\maketitle

\begin{abstract}
This paper presents a hybrid information retrieval system for scientific articles on artificial intelligence applied to public procurement. The system combines BM25 (sparse retrieval) with a TF-IDF/KNN-based dense retriever via Reciprocal Rank Fusion (RRF) over a bilingual corpus (English + Portuguese). We collected 3,106 articles from OpenAlex, BDTD, and arXiv, created 15 bilingual queries with manual relevance judgments, and evaluated performance using P@10, R@10, MAP, and nDCG@10. Results show the supervised re-ranking module (M1) achieves the best MAP (0.864), with BM25 a strong baseline (MAP=0.830). The dense retriever scored lower, which we trace to pooling bias in the relevance judgments rather than to the model itself. We discuss how this annotation bottleneck motivates semi-supervised learning for the procurement domain.
\end{abstract}

\begin{resumo}
Este trabalho apresenta um sistema de recuperação de informação híbrido para artigos científicos sobre inteligência artificial aplicada a compras públicas. O sistema combina BM25 (recuperação esparsa) com um recuperador denso baseado em TF-IDF/KNN via Reciprocal Rank Fusion (RRF) sobre um corpus bilíngue (inglês + português). Foram coletados 3.106 artigos da OpenAlex, da BDTD e do arXiv, criadas 15 consultas bilíngues com julgamentos manuais de relevância, e avaliado o desempenho com P@10, R@10, MAP e nDCG@10. Os resultados mostram o módulo de re-ranking supervisionado (M1) com o melhor MAP (0,864), e o BM25 como baseline forte (MAP=0,830). O recuperador denso pontuou abaixo, o que atribuímos ao viés de pooling nos julgamentos de relevância, e não ao modelo em si. Discutimos como esse gargalo de anotação motiva o aprendizado semissupervisionado para o domínio.
\end{resumo}

\section{Introdução}

O volume de publicações científicas sobre inteligência artificial aplicada a domínios governamentais cresce rapidamente. Pesquisadores e servidores públicos que buscam referências para projetos de automação em compras públicas enfrentam a dificuldade de localizar artigos relevantes em bases amplas e heterogêneas. Sistemas de recuperação de informação (IR) específicos para este domínio ainda são escassos. Os sistemas de Geração Aumentada por Recuperação (RAG) \cite{lewis2020,gao2024} atacam esse problema: acoplam um recuperador a um modelo generativo. Mas a qualidade da resposta final depende do recuperador. É o ``R'' do RAG, e é o objeto deste trabalho.

Este trabalho implementa e avalia um sistema de recuperação de artigos científicos sobre IA aplicada a compras públicas e licitações, comparando três abordagens: (i) BM25 como baseline esparso, (ii) KNN com representação TF-IDF como recuperador denso, e (iii) um módulo híbrido que combina ambos via Reciprocal Rank Fusion (RRF).

A contribuição principal é empírica. Avaliamos três recuperadores (BM25, KNN/TF-IDF e denso), a fusão híbrida e os módulos de aprofundamento sobre uma coleção bilíngue de 3.106 artigos do domínio, com 15 consultas e julgamentos manuais de relevância. A comparação mostra os trade-offs entre recuperação esparsa, densa e híbrida, e expõe o gargalo da anotação manual de relevância em um domínio nichado e bilíngue.

\section{Trabalhos Relacionados}

\textbf{BM25 e o paradigma probabilístico.} O BM25 \cite{robertson2009,manning2008} é derivado do framework de relevância probabilística, sendo uma extensão prática do modelo binário de independência (relacionado ao Naïve Bayes). A função de pontuação incorpora frequência do termo (TF), frequência inversa de documento (IDF) e normalização por comprimento do documento, parametrizada por $k_1$ (saturação de TF) e $b$ (impacto do comprimento).

\textbf{Recuperação densa e fusão.} Karpukhin et al. \cite{karpukhin2020} mostraram, com o Dense Passage Retrieval, que representações densas aprendidas superam o BM25 em perguntas e respostas, embora o BM25 siga competitivo em domínios de vocabulário controlado \cite{thakur2021}. Cormack et al. \cite{cormack2009} demonstraram que o Reciprocal Rank Fusion supera métodos individuais sem treino adicional. Lin \cite{lin2021} unifica recuperação esparsa, densa e híbrida sob uma visão representacional, o que fundamenta a comparação experimental deste trabalho.

\textbf{NLP para domínio legal.} Trabalhos recentes exploram modelos de linguagem para textos jurídicos \cite{chalkidis2020}, mas poucos abordam especificamente o domínio de compras públicas e licitações brasileiras.

\section{Metodologia}

\subsection{Escopo da coleção e relação com o projeto de pesquisa}

Este trabalho é o embrião técnico de uma proposta de mestrado, ainda em fase de ideia a validar com os orientadores. A proposta trata da classificação semissupervisionada de textos licitatórios para apoio à análise de conformidade em compras públicas. O sistema de recuperação aqui desenvolvido vira, nessa proposta, o módulo de busca por peças processuais similares. Construí um corpus de artigos científicos para validar a arquitetura antes de migrar para os documentos reais (pareceres, editais e atas), que dependem de coleta via LAI e de anotação própria.

A coleção foi construída a partir da API OpenAlex em duas etapas: (i) coleta de artigos em inglês sobre IA aplicada a procurement e legal NLP; (ii) coleta complementar de artigos em português sobre licitações, direito administrativo e auditoria governamental. A inclusão de documentos em português é motivada pelo fato de que o domínio-alvo (compras públicas brasileiras, Lei 14.133/2021) produz literatura predominantemente em português.

A segunda etapa exigiu uma busca prévia: quais bases e periódicos brasileiros publicam sobre o cruzamento entre tecnologia e administração pública? Conduzimos essa busca com apoio do Gemini (Deep Research) e verificamos os resultados manualmente, o que definiu as palavras-chave em português. A coleta usou a API OpenAlex, que indexa periódicos brasileiros e dá, numa única interface, acesso amplo à produção nacional sobre o tema. Foi a escolha certa para esta etapa: gratuita, sem rate limit agressivo e com boa cobertura do domínio.

Quanto à conformidade com a Seção 5 do enunciado: a OpenAlex e o arXiv constam entre as fontes explicitamente permitidas, e a BDTD (Biblioteca Digital Brasileira de Teses e Dissertações, do Ibict) foi adotada como fonte equivalente para a cobertura em português, mediante aprovação do professor. Nenhuma coleção pré-pronta de \emph{benchmark} de IR (e.g., BEIR, MS MARCO) foi usada como coleção principal; a coleção foi inteiramente construída e documentada para este trabalho.

Incluímos também o arXiv, e o resultado foi revelador. Buscas por frase exata sobre o tema retornaram apenas 42 artigos on-topic. O arXiv é forte em aprendizado de máquina, mas raso em compras públicas e direito administrativo brasileiro. Esse vazio não é uma falha da coleta. É a própria lacuna que motiva a pesquisa de mestrado: há muita IA e muita licitação, mas pouco no cruzamento dos dois.

Uma nota sobre o formato dos dados. O campo identificador de cada documento chama-se \texttt{arxiv\_id}. O nome é herança do material de apoio do professor (o notebook de coleta do arXiv). Mantivemos o nome por compatibilidade, mas ele guarda o id da fonte de origem: id do arXiv para artigos do arXiv, id da OpenAlex (prefixo \texttt{W}) para a OpenAlex e id da BDTD para as teses. Um campo \texttt{source} registra a proveniência de cada documento.

As bases governamentais mapeadas (PNCP, Compras.gov.br, TCU, LexML) e os pareceres obtidos via LAI não entram nesta coleção. Elas são o corpus-alvo da etapa de mestrado, quando o foco deixa de ser artigos científicos e passa a ser peças processuais reais. Esses documentos têm uma natureza própria. O PNCP e os sistemas de compras entregam cada processo com um número variável de artefatos, que muda conforme o trâmite. Os resumos e textos introdutórios não seguem padronização real, o que abre espaço para inferências erradas na automação. É aí que o aprendizado semissupervisionado faz sentido: rotular tudo à mão é inviável pelo tamanho do conjunto, então parte-se de poucos rótulos e propaga-se o aprendizado para o resto.

\begin{itemize}
    \item \textbf{Fontes:} OpenAlex (artigos de periódicos), BDTD/Ibict (teses e dissertações brasileiras) e arXiv (pré-prints)
    \item \textbf{Palavras-chave (EN):} ``public procurement NLP'', ``legal document classification'', ``regulatory compliance automation'', entre outras
    \item \textbf{Palavras-chave (PT):} ``licitação inteligência artificial'', ``tribunal contas auditoria'', ``processamento linguagem natural documentos jurídicos'', ``contratação pública tecnologia'', entre outras
    \item \textbf{Janela temporal:} 2018--2026
    \item \textbf{Tamanho final:} 3.106 documentos (2.758 OpenAlex + 306 BDTD + 42 arXiv; 2.135 EN + 971 PT) após deduplicação
\end{itemize}

A distribuição temporal é equilibrada e o corpus cobre múltiplas subáreas: IA em direito, compliance, procurement, auditoria governamental e processamento de textos jurídicos.

\subsection{Pré-processamento}

O mesmo pipeline é aplicado a documentos e consultas, com detecção automática de idioma para seleção do stemmer adequado:

\begin{enumerate}
    \item \textbf{Lowercase:} conversão para minúsculas.
    \item \textbf{Tokenização:} expressão regular que mantém palavras alfabéticas com mais de 1 caractere.
    \item \textbf{Stopwords:} lista NLTK bilíngue (inglês + português) expandida com termos genéricos de abstracts em ambos os idiomas.
    \item \textbf{Detecção de idioma:} heurística baseada em palavras-marcador do português (``de'', ``da'', ``para'', ``são'', etc.). Se 3 ou mais marcadores aparecem nas primeiras 50 palavras, o texto é classificado como português.
    \item \textbf{Stemming:} Porter Stemmer para inglês; RSLP Stemmer para português (ambos do NLTK).
\end{enumerate}

O texto indexado de cada documento é a concatenação de título e abstract. A abordagem bilíngue funciona assim: cada texto passa pelo detector de idioma e recebe o stemmer da sua língua (Porter para inglês, RSLP para português). As stopwords dos dois idiomas são removidas sempre. Com isso, BM25 e TF-IDF tratam PT e EN no mesmo índice, sem separar as subcoleções. O alcance, porém, é limitado. Uma consulta em português só recupera um documento em inglês quando os dois compartilham radicais após o stemming, o que costuma ocorrer só com nomes próprios e siglas. Para casar significado entre idiomas sem coincidência de termos, recorremos ao recuperador denso (Seção~\ref{sec:dense}).

\subsection{BM25 (recuperação esparsa)}

Usamos a implementação \texttt{rank\_bm25} (BM25Okapi) com os hiperparâmetros padrão da literatura: $k_1=1.5$ e $b=0.75$ \cite{robertson2009}. A escolha segue Robertson e Zaragoza. Eles mostram que esses valores são robustos para coleções de tamanho moderado, sem precisar de ajuste empírico.

A pontuação BM25 para um documento $d$ dada uma consulta $q$ com termos $q_1, \ldots, q_n$ é:

\begin{equation}
\text{score}(d, q) = \sum_{i=1}^{n} \text{IDF}(q_i) \cdot \frac{f(q_i, d) \cdot (k_1 + 1)}{f(q_i, d) + k_1 \cdot (1 - b + b \cdot \frac{|d|}{\text{avgdl}})}
\end{equation}

\textbf{Conexão com Naïve Bayes.} BM25 e Naïve Bayes partem da mesma pergunta: qual a chance de um documento ser relevante, dados os termos da consulta? O Naïve Bayes responde assumindo que os termos são independentes entre si. O BM25 herda essa ideia. A componente IDF é o coração da ligação. Ela pesa cada termo pelo quanto ele distingue um documento relevante dos demais, o mesmo papel do log-odds que o Naïve Bayes calcula para cada palavra dada a classe. Termos raros valem mais porque carregam mais informação. Sobre essa base, o BM25 acrescenta dois ajustes que o Naïve Bayes puro não tem. O $k_1$ controla a saturação: a décima ocorrência de um termo agrega menos que a segunda, porque repetir não aumenta a relevância para sempre. O $b$ corrige o tamanho do documento, para que um texto longo não pareça mais relevante só por conter mais palavras. Em resumo, o BM25 é um Naïve Bayes calibrado para ranquear.

\subsection{KNN/TF-IDF (recuperação por vizinhos mais próximos)}

\textbf{O que é TF-IDF.} O TF-IDF transforma cada documento em um vetor numérico. Cada posição do vetor é um termo do vocabulário. O valor combina dois fatores: TF, quantas vezes o termo aparece no documento, e IDF, quão raro o termo é no corpus inteiro. Um termo frequente naquele documento e raro nos demais recebe peso alto, porque é o que o caracteriza. Termos comuns a quase todos os textos recebem peso baixo. Documentos sobre o mesmo assunto terminam com valores parecidos nas mesmas posições, e seus vetores ficam próximos no espaço.

Construímos a matriz TF-IDF com o \texttt{TfidfVectorizer} do \texttt{scikit-learn}. Os cortes \texttt{min\_df=2} e \texttt{max\_df=0.95} descartam termos que aparecem em um só documento e termos presentes em mais de 95\% deles, que não ajudam a distinguir nada. A busca devolve os $K$ documentos de maior similaridade de cosseno com a consulta.

\textbf{Conexão com a disciplina:} O uso de KNN aqui difere da classificação vista em aula em um aspecto conceitual: em vez de encontrar os K vizinhos mais próximos e votar uma classe, devolvemos os K documentos mais similares como ranking de resposta à consulta. O algoritmo é o mesmo (busca por proximidade em espaço vetorial), mas a saída é uma lista ranqueada e não um rótulo. A ponderação TF-IDF conecta com a discussão de Naïve Bayes: o componente IDF é análogo à probabilidade a priori de observar cada termo, penalizando termos frequentes e valorizando termos discriminativos.

\textbf{Escolha da similaridade:} Optamos pela similaridade do cosseno em vez da distância euclidiana porque o cosseno normaliza pelo comprimento dos vetores, tornando a métrica invariante ao tamanho do documento. Dois abstracts sobre o mesmo tema, um curto e um longo, terão vetores na mesma direção (alto cosseno) mesmo com magnitudes diferentes (alta distância euclidiana).

\subsection{Recuperador denso com embeddings}
\label{sec:dense}

O TF-IDF casa termos exatos. Ele não sabe que ``licitação'' e ``procurement'' falam da mesma coisa. Para cobrir esse caso, adicionamos um terceiro recuperador, baseado em embeddings. Um embedding é um vetor denso produzido por uma rede neural pré-treinada, no qual textos de significado parecido ficam próximos mesmo sem compartilhar palavras. Usamos o modelo multilíngue \texttt{paraphrase-multilingual-MiniLM-L12-v2} \cite{reimers2019}, que mapeia português e inglês no mesmo espaço. Cada documento e cada consulta viram um vetor de 384 dimensões, e a busca devolve os $K$ vizinhos de maior cosseno. O uso de modelos prontos é permitido (a aula não cobre o treinamento deles), então implementamos apenas a codificação e a busca. A motivação é direta: resolver a recuperação cross-lingual que o TF-IDF não alcança.\footnote{Hardware: Intel Core i7-12650H (16 threads), Debian 13, Python 3.13, execução em CPU (sem GPU). A codificação dos 3.106 documentos levou cerca de 3,6 minutos; BM25 e TF-IDF rodam em segundos.}

\subsection{Módulos de aprofundamento}

Além do pipeline obrigatório, implementamos os cinco módulos da Seção 4 do enunciado. O Mestrado exige só um. Fizemos os cinco por interesse próprio em comparar as abordagens. Cada módulo é avaliado na Seção~\ref{sec:resultados}.

\textbf{M5: ranking híbrido (RRF).} O RRF \cite{cormack2009} combina rankings sem normalizar scores:

\begin{equation}
\text{score}_{\text{RRF}}(d) = \sum_{r \in R} \frac{1}{k + \text{rank}_r(d)}
\end{equation}

onde $k=60$ (constante padrão) e $R$ reúne os rankings do BM25 e do recuperador denso. A vantagem é não depender de calibração entre escalas: o BM25 varia de 0 a cerca de 20 e o cosseno de 0 a 1, mas o RRF olha só a posição. É um ensemble. Cada sistema vota pelo lugar que deu ao documento, e juntar votos de modelos com vieses diferentes reduz o erro, como na combinação de classificadores.

\textbf{M1: re-ranking supervisionado.} Treinamos uma Regressão Logística para reordenar o top-100 do BM25. As features de cada par (consulta, documento) são cinco: score BM25, cosseno TF-IDF, similaridade densa, sobreposição de termos e ano. O rótulo vem dos qrels. Para não treinar e testar na mesma consulta, usamos validação leave-one-query-out: o modelo treina em 14 consultas e reordena a 15ª. Liga à aula de classificação supervisionada.

\textbf{M2: agrupamento dos resultados.} Aplicamos K-means sobre o top-30 de uma consulta, no espaço dos embeddings. O número de grupos é escolhido pela silhueta. Cada grupo é descrito pelos seus termos mais frequentes, formando facetas temáticas. Liga às aulas de agrupamento de dados.

\textbf{M3: expansão de consulta por regras de associação.} Tratamos cada documento como uma transação de termos. Para cada termo da consulta, mineramos regras $\{A\} \Rightarrow \{B\}$ com suporte, confiança e lift acima de um limiar. O lift descarta termos genéricos, que coocorrem com quase tudo. Os termos consequentes entram na consulta, que é re-executada no BM25. Liga à aula de regras de associação.

\textbf{M4: otimização de hiperparâmetros.} Fizemos grid search sobre $k_1 \in \{1{,}2;\ 1{,}5;\ 2{,}0\}$ e $b \in \{0{,}5;\ 0{,}75;\ 1{,}0\}$ do BM25, maximizando o MAP no conjunto de avaliação. A configuração vencedora e a tabela completa de desempenho constam na Seção~\ref{sec:resultados}. Liga à aula de algoritmos de otimização.

\subsection{Avaliação}

\subsubsection{Consultas de teste}

Criamos 15 consultas representativas do domínio (Tabela~\ref{tab:queries}). Elas não são aleatórias. Espelham as subáreas da proposta de mestrado: classificação de documentos (q01, q15), conformidade e compliance (q02, q08, q11), análise de editais e contratos (q03, q04, q13), fraude e auditoria (q05, q12), apoio à decisão, avaliação e grafos (q06, q07, q09) e o recorte central da pesquisa, a classificação semissupervisionada com poucos rótulos (q10). Variamos a especificidade de propósito: algumas amplas, outras finas.

A divisão entre 10 consultas em inglês e 5 em português também é deliberada. A literatura de método (NLP, IR, ML) é quase toda em inglês, e é nela que estão os trabalhos técnicos. As cinco consultas em português miram o domínio brasileiro: Lei 14.133, tribunal de contas, contratos públicos. A divisão testa o sistema nos dois eixos que o problema real combina: a técnica, em inglês, e a aplicação, em português.

\begin{table}[h]
\centering
\caption{Consultas de teste usadas na avaliação (10 EN + 5 PT)}
\label{tab:queries}
\small
\begin{tabular}{lll}
\toprule
\textbf{ID} & \textbf{Idioma} & \textbf{Consulta} \\
\midrule
q01 & EN & public procurement document classification machine learning \\
q02 & EN & legal compliance checking NLP automated verification \\
q03 & EN & tender document analysis natural language processing \\
q04 & EN & government contract clause extraction deep learning \\
q05 & EN & procurement fraud detection anomaly bidding collusion \\
q06 & EN & e-procurement automation decision support system \\
q07 & EN & bid evaluation criteria similarity neural networks \\
q08 & EN & regulatory compliance verification document NLP \\
q09 & EN & knowledge graph public procurement legislation \\
q10 & EN & semi-supervised classification legal text few labels \\
q11 & PT & licitação análise conformidade inteligência artificial \\
q12 & PT & tribunal contas auditoria automatizada documentos \\
q13 & PT & contratos públicos processamento linguagem natural \\
q14 & PT & lei 14133 compras públicas tecnologia \\
q15 & PT & classificação documentos administrativos aprendizado máquina \\
\bottomrule
\end{tabular}
\end{table}

\subsubsection{Construção dos qrels}

Anotamos a relevância por pooling. Para cada consulta, unimos o top-10 de todos os sistemas (BM25, KNN, denso, RRF e os módulos), removemos duplicatas e julgamos cada documento na escala 0/1/2. O conjunto começou com 231 julgamentos manuais. Ao acrescentar o denso, o arXiv e os módulos, o pool cresceu e exigiu 212 julgamentos novos. Esses foram anotados com apoio de um classificador treinado nos 231 originais e revisados pelo autor, totalizando 443 julgamentos. A Seção~\ref{sec:anotacao} discute o que essa etapa revelou e por que ela é central para a pesquisa de mestrado.

\textbf{Critério de relevância adotado:} consideramos relevante (1) todo artigo que aborda diretamente a técnica ou domínio especificado na query, e altamente relevante (2) artigos que tratam exatamente do cruzamento entre a técnica de NLP/ML indicada e o domínio de procurement/legal da query.

\textbf{Métricas reportadas:} Precision@10 (P@10), Recall@10 (R@10), Mean Average Precision (MAP) e normalized Discounted Cumulative Gain (nDCG@10).

\section{Resultados}
\label{sec:resultados}

\begin{table}[h]
\centering
\caption{Resultados médios sobre 15 consultas (corpus de 3.106 documentos, 443 julgamentos). M2 e M4 não geram ranking comparável e são reportados à parte.}
\label{tab:results}
\begin{tabular}{lcccc}
\toprule
\textbf{Sistema} & \textbf{P@10} & \textbf{R@10} & \textbf{MAP} & \textbf{nDCG@10} \\
\midrule
BM25                & 0.727 & 0.644 & 0.830 & 0.791 \\
KNN/TF-IDF          & 0.720 & 0.639 & 0.740 & 0.697 \\
Denso (embeddings)  & 0.240 & 0.172 & 0.180 & 0.306 \\
RRF (BM25+denso)    & 0.580 & 0.456 & 0.552 & 0.642 \\
M1 (re-ranking)     & \textbf{0.740} & \textbf{0.651} & \textbf{0.864} & \textbf{0.794} \\
M3 (expansão)       & 0.567 & 0.515 & 0.608 & 0.643 \\
\bottomrule
\end{tabular}
\end{table}

O M1 (re-ranking supervisionado) é o melhor sistema: MAP 0,864 e P@10 0,740, acima do BM25 (MAP 0,830). Faz sentido, porque o M1 aprende a combinar os sinais dos outros recuperadores. O BM25 é um baseline forte, e o KNN/TF-IDF empata com ele em P@10. O recuperador denso e o RRF aparecem bem abaixo. Esse número pede cuidado na leitura. Ele diz mais sobre como construímos o gabarito do que sobre a qualidade do denso, e a Seção~\ref{sec:anotacao} explica por quê.

\subsection{Análise qualitativa}

\textbf{O re-ranking em ação (q01).} Na consulta ``public procurement document classification machine learning'', o M1 melhora o BM25 em toda a linha: P@10 de 0,90 para 1,00, MAP de 0,78 para 0,89 e nDCG de 0,90 para 0,96. O classificador aprende a combinar score BM25, cosseno, similaridade densa e sobreposição de termos, e empurra para o topo os documentos certos. É o ganho que se espera de um re-ranking supervisionado.

\textbf{O casamento exato ainda resolve muito (q14).} A consulta em português ``lei 14133 compras públicas tecnologia'' tem P@10=1,00 no BM25. Quando a consulta traz o termo exato do domínio, como o número da lei, o esparso acerta sem precisar de semântica.

\textbf{O denso e o limite do gabarito (q07).} O recuperador denso pontua baixo em quase todas as consultas, chegando a P@10=0,00 em q07. Mas, ao ler os resultados, vê-se que ele recupera documentos plausíveis. Em q01, ele traz a dissertação ``Classificação automática de documentos'' (similaridade 0,64), que o nosso gabarito marcou como não-relevante. O problema não está no denso. Está no gabarito: ele foi montado a partir de um pool majoritariamente lexical, e os documentos que só o denso encontra ficaram de fora ou foram subjulgados. A próxima seção trata disso.

\textbf{Módulos M2 e M4.} O M2 (agrupamento) separou o top-30 de uma consulta em grupos coerentes, escolhidos por silhueta, com facetas temáticas distintas. O M4 (otimização) varreu a grade de $k_1$ e $b$ do BM25 e encontrou a melhor configuração no conjunto de avaliação. As saídas completas estão nos apêndices e no repositório.

\section{Discussão}

\subsection{Trade-offs entre recuperação esparsa e densa}

\textbf{O melhor sistema foi o M1.} O re-ranking supervisionado lidera em todas as métricas (MAP 0,864, P@10 0,740), porque combina os sinais dos demais recuperadores em vez de depender de um só. O BM25 é o baseline forte (MAP 0,830), e o KNN/TF-IDF empata com ele em P@10 (0,720 contra 0,727). A subcoleção em português, mais diversa, ajuda o TF-IDF a competir com o esparso.

\textbf{O denso e o RRF ficaram abaixo, e isso merece explicação.} Pelos números, o denso (P@10 0,240) e o RRF (0,580) decepcionam. Mas a próxima seção mostra que o problema está no gabarito, não nos modelos. A barreira de idioma também pesa: o stemmer bilíngue casa pouco entre PT e EN, e era o denso multilíngue que deveria cobrir esse vão. Para medir esse ganho, porém, é preciso um gabarito que enxergue os relevantes que o denso encontra.

\subsection{O desafio da anotação de relevância}
\label{sec:anotacao}

Ao adicionar o denso, o arXiv e os módulos, o pool de avaliação cresceu. Só o denso trouxe 120 documentos novos ao top-10. Anotar tudo à mão, relendo cada par consulta-documento, é caro. Recorremos a um classificador treinado nos 231 julgamentos originais para sugerir os 212 rótulos novos, e os revisamos. Esse classificador aprendeu com um gabarito de origem lexical, então tende a marcar como não-relevante justamente os documentos que o denso encontra por significado. O resultado é que o denso é penalizado por achar itens certos que o gabarito não reconhece.

O fenômeno é conhecido. Zobel \cite{zobel1998} mostrou que documentos fora do pool, contados como não-relevantes, enviesam a comparação contra sistemas que recuperam itens novos. Voorhees \cite{voorhees2000} mostrou que julgamentos de relevância variam entre anotadores, mas ainda permitem comparar sistemas. Sanderson \cite{sanderson2010} descreve como construir coleções de teste mais sólidas. A relevância graduada que adotamos segue Järvelin e Kekäläinen \cite{jarvelin2002}.

Apresentamos os números como eles são. Eles refletem o método e o conhecimento que tínhamos no momento, não um veredito final sobre o denso. O caminho de melhoria é claro: pool mais profundo, mais de um anotador com acordo medido, e julgamento por leitura, não por proxy lexical.

E aqui está o elo com a pesquisa de mestrado. Anotar relevância à mão é lento, subjetivo e não escala. Sentimos isso ao construir 443 julgamentos para apenas 15 consultas. No domínio licitatório, com volume muito maior e sem dados rotulados, o custo se torna proibitivo. É essa a motivação do mestrado: aprendizado semissupervisionado para classificar e recuperar com poucos rótulos, reduzindo a dependência da anotação manual.

\subsection{Conexões com a disciplina}

Os componentes mapeiam diretamente para tópicos vistos em sala. BM25 e TF-IDF ligam ao paradigma probabilístico do Naïve Bayes (IDF como log-odds). O recuperador denso e o KNN ligam à busca por vizinhos em espaço vetorial; a diferença para a aula é a saída, um ranking em vez de um rótulo. O RRF liga ao ensemble de classificadores. E os módulos M1 a M4 ligam, na ordem, à classificação supervisionada, ao agrupamento de dados, às regras de associação e aos algoritmos de otimização.

\subsection{Limitações}

(i) O gabarito tem viés de pooling, e parte dos julgamentos novos foi sugerida por classificador e revisada, não anotada do zero (Seção~\ref{sec:anotacao}). (ii) A avaliação usa 15 consultas e um único anotador, sem acordo inter-anotador; o M4 otimiza sobre esse mesmo conjunto, com risco de superajuste. (iii) O corpus mistura subáreas adjacentes (blockchain, ética de IA), o que introduz ruído. (iv) O recuperador denso usa um modelo pré-treinado genérico; um modelo jurídico em português tende a ir melhor.

\section{Conclusão}

Implementamos e avaliamos um sistema de recuperação híbrido para artigos sobre IA em compras públicas. O corpus é bilíngue, com 3.106 documentos: artigos da OpenAlex (inglês e português), teses e dissertações da BDTD e pré-prints do arXiv. O melhor sistema foi o re-ranking supervisionado (M1), com MAP 0,864. Mais importante que o ranking, porém, foi uma lição de método. O recuperador denso ficou subavaliado porque o gabarito, de origem lexical, não reconhece os relevantes que ele encontra por significado. Isso expõe o gargalo da anotação manual de relevância. Como trabalho futuro, pretendemos (i) reconstruir o gabarito com pool mais profundo e mais de um anotador, para medir o denso de forma justa; (ii) incorporar dados primários das APIs governamentais (PNCP, Compras.gov.br, TCU, LexML) e pareceres via LAI; e (iii) avançar para a classificação semissupervisionada de textos licitatórios, foco do mestrado\footnote{Vídeo de apresentação e repositório: ver arquivo LINKS.txt.}.

Uma extensão conceitual está em prototipação, em ramo separado do repositório (\texttt{v2-reference-finder}), sem impacto sobre os resultados aqui reportados. Ela reaproveita o núcleo de recuperação como assistente de referências bibliográficas. O texto de um artigo em elaboração entra como consulta longa. Dele extraímos palavras-chave que alimentam a busca híbrida. A saída são referências em formato BibTeX e uma análise de lacunas: temas presentes no texto, mas sem cobertura no corpus. O recorte transpõe o problema de recuperação deste trabalho para o apoio à revisão de literatura, em continuidade com a pesquisa de mestrado.

Para fins de reprodução e avaliação, o sistema completo (pipeline, dados e relatório) está empacotado em container e disponível para teste interativo por um terminal web no navegador\footnote{O endereço de acesso (terminal web, sem necessidade de instalar nada) está no arquivo LINKS.txt da submissão. O acesso conduz diretamente ao ambiente isolado do trabalho, onde \texttt{make eval} reproduz as métricas aqui reportadas e \texttt{make demo} executa uma consulta. Um segundo endereço dá acesso ao protótipo \texttt{v2-reference-finder}.}.

\subsection*{Declaração de uso de IA generativa}

Conforme Resolução nº 455-COUN/UFMS (BO-UFMS ID 581705; Art. 6º, V e Art. 12, Parágrafo único), declara-se o uso das seguintes ferramentas de IA como apoio durante a elaboração deste trabalho:

\begin{itemize}
    \item \textbf{Gemini (Google, funcionalidade Deep Research):} usado para pesquisa exploratória de fontes e bases de dados acadêmicas em língua portuguesa. A especificidade do domínio (compras públicas brasileiras, Lei 14.133/2021) e a escassez de artigos no ArXiv sobre o tema motivaram a busca por bases complementares. O Gemini auxiliou na identificação de periódicos e repositórios brasileiros relevantes, cujos resultados foram avaliados criticamente e verificados manualmente pelo autor.
    \item \textbf{Claude/Kiro (Anthropic):} apoio na implementação de scripts de coleta e pipeline, sugestões de hiperparâmetros e estruturação do relatório.
\end{itemize}

Em ambos os casos, as ferramentas serviram como apoio instrumental. As decisões de projeto (escolha do tema, definição de queries, critérios de relevância), a anotação manual dos qrels, a análise crítica dos resultados e a redação argumentativa do relatório são integralmente do autor. Todo código gerado com auxílio de IA foi revisado, testado e adaptado antes da incorporação.

\bibliographystyle{sbc}
\bibliography{references}

\section*{Apêndice A: Average Precision por consulta}

A Tabela~\ref{tab:apx} traz a AP de cada sistema por consulta. Ela mostra de forma transparente o efeito discutido na Seção~\ref{sec:anotacao}: o denso pontua perto de zero em consultas como q12 e q15, onde os documentos que ele recupera não constam como relevantes no gabarito.

\begin{table}[h]
\centering
\small
\caption{Average Precision (AP) por consulta.}
\label{tab:apx}
\begin{tabular}{lcccccc}
\toprule
\textbf{qid} & \textbf{BM25} & \textbf{KNN} & \textbf{Denso} & \textbf{RRF} & \textbf{M1} & \textbf{M3} \\
\midrule
q01 & 0.782 & 0.693 & 0.272 & 0.630 & 0.891 & 0.647 \\
q02 & 0.590 & 0.738 & 0.189 & 0.448 & 0.734 & 0.571 \\
q03 & 0.759 & 0.692 & 0.294 & 0.683 & 0.824 & 0.478 \\
q04 & 0.813 & 0.676 & 0.349 & 0.702 & 0.855 & 0.484 \\
q05 & 0.850 & 0.758 & 0.203 & 0.570 & 0.909 & 0.401 \\
q06 & 0.750 & 0.423 & 0.200 & 0.606 & 0.759 & 0.344 \\
q07 & 0.729 & 0.436 & 0.028 & 0.392 & 0.714 & 0.282 \\
q08 & 0.788 & 0.720 & 0.301 & 0.661 & 0.886 & 0.525 \\
q09 & 0.927 & 0.895 & 0.190 & 0.555 & 0.977 & 0.445 \\
q10 & 0.713 & 0.805 & 0.154 & 0.614 & 0.795 & 0.753 \\
q11 & 1.000 & 0.982 & 0.012 & 0.595 & 1.000 & 0.810 \\
q12 & 0.952 & 0.798 & 0.001 & 0.433 & 0.954 & 0.864 \\
q13 & 1.000 & 0.700 & 0.250 & 0.532 & 0.833 & 1.000 \\
q14 & 0.959 & 0.919 & 0.253 & 0.705 & 0.987 & 0.672 \\
q15 & 0.841 & 0.861 & 0.000 & 0.149 & 0.841 & 0.841 \\
\bottomrule
\end{tabular}
\end{table}

\section*{Apêndice B: Otimização do BM25 (M4)}

Grid search de $k_1$ e $b$, com MAP no conjunto de avaliação. A melhor configuração ($k_1=2{,}0$, $b=0{,}5$) eleva o MAP de 0,830 (padrão) para 0,851.

\begin{table}[h]
\centering
\small
\caption{MAP por configuração do BM25.}
\begin{tabular}{ccc|ccc|ccc}
\toprule
$k_1$ & $b$ & MAP & $k_1$ & $b$ & MAP & $k_1$ & $b$ & MAP \\
\midrule
1.2 & 0.50 & 0.817 & 1.5 & 0.50 & 0.829 & 2.0 & 0.50 & \textbf{0.851} \\
1.2 & 0.75 & 0.816 & 1.5 & 0.75 & 0.830 & 2.0 & 0.75 & 0.831 \\
1.2 & 1.00 & 0.781 & 1.5 & 1.00 & 0.774 & 2.0 & 1.00 & 0.782 \\
\bottomrule
\end{tabular}
\end{table}

\section*{Apêndice C: Agrupamento dos resultados (M2)}

Para a consulta ``lei 14133 compras públicas tecnologia'', o K-means sobre o top-30 (espaço dos embeddings) escolheu $k=2$ pela silhueta. As facetas, descritas pelos termos mais frequentes de cada grupo:

\begin{itemize}
    \item \textbf{Grupo 0 (marco legal e processo):} público, lei, compra, processo, contrato, licitação.
    \item \textbf{Grupo 1 (tecnologia e sustentabilidade):} compra, público, análise, tecnologia, sustentável, federal.
\end{itemize}

O pool completo de anotação, com sugestões, sinais e julgamentos, está disponível como planilha no repositório\footnote{\url{https://github.com/luizfpq/ufms-facom-ia-2026.1-t1/blob/main/eval/pool_revisao.xlsx}}.

\end{document}
